
\documentclass[a4paper,12pt]{article}
\usepackage{amsmath, amssymb}
\usepackage{xcolor}
\usepackage{tcolorbox}
\usepackage{multicol}
\usepackage{geometry}
\hyphenpenalty=100000

\geometry{left=1.5cm, right=1.5cm, top=2cm, bottom=2cm}

\tcbuselibrary{theorems}

\definecolor{PurpleEx}{HTML}{5d478b}
\definecolor{GrayEx}{HTML}{f2f2f2}
\definecolor{GrayBG}{HTML}{d9d9d9}
\definecolor{BlackSav}{HTML}{000000}

\title{
\centering
\tcbset{colframe=BlackSav,colback=GrayBG}
\tcbox[tikznode]{
Chapitre 3 \\ Nombres complexes - propriétés géométriques
}
}
\usepackage{tikz}
\author{}
\date{}

\newtcbtheorem
  []% init options
  {exercice}% name
  {Exercice}% title
  {%
    colback=GrayEx,
    colframe=PurpleEx,
    fonttitle=\bfseries,
  }% options
  {ex}% prefix


\begin{document}

\maketitle

\begin{multicols}{2}

\begin{exercice}[title=Exercice 1 \quad$\star$\quad\ {[}Calculer{]}]{}{}
Calculer le module et l’argument des nombres complexes suivants :
\begin{enumerate}
    \item \( z = 1 + i \)
    \item \( z = 2\pi i \)
    \item \( z = 2 - 2\sqrt{3}i \)
    \item \( z = -5 - 5i \)
    \item \( z = \sqrt{6} - \sqrt{2}i \)
\end{enumerate}
\end{exercice}

\begin{exercice}[title=Exercice 2 \quad$\star$\quad\ {[}Calculer{]}]{}{}
Déterminer la forme trigonométrique et la forme exponentielle des nombres complexes suivants :
\begin{enumerate}
    \item \( z = 2 + 2i \)
    \item \( z = -3i \)
    \item \( z = 1 + \sqrt{3}i \)
    \item \( z = \frac{\sqrt{3}}{3} - \frac{1}{3}i \)
    \item \( z = -\sqrt{5} - \sqrt{5}i \)
\end{enumerate}
\end{exercice}

\begin{exercice}[title=Exercice 3 \quad$\star$\quad\ {[}Calculer{]}]{}{}
Déterminer la forme algébrique de :
\[
(-1 + i)^{13}.
\]
\end{exercice}

\begin{exercice}[title=Exercice 4 \quad$\star\star$\quad\ {[}Calculer{]}]{}{}
Déterminer la forme exponentielle de :
\[
e^{i\frac{\pi}{4}} + e^{i\frac{\pi}{3}}.
\]
\end{exercice}

\begin{exercice}[title=Exercice 5 \quad$\star\star$\quad\ {[}Calculer{]}]{}{}
Calculer
\[
\int_{\frac{2\pi}{3}}^{0} \cos^3(t) \, dt.
\]
\end{exercice}

\begin{exercice}[title=Exercice 6 \quad$\star\star$\quad\ {[}Calculer{]}]{}{}
Calculer
\[
\int_{\frac{\pi}{4}}^{0} \cos^2(t) \sin^2(t) \, dt.
\]
\end{exercice}

\begin{exercice}[title=Exercice 7 \quad$\star\star$\quad\ {[}Calculer{]}]{}{}
À l’aide de la formule de Moivre, exprimer \( \cos(3x) \) et \( \sin(3x) \) en fonction de \( \cos(x) \) et \( \sin(x) \).
\end{exercice}

\begin{exercice}[title={Exercice 8 \quad$\star\star$\quad\ [Raisonner]}]{}{}
Montrer que pour tout \( z \in \mathbb{U} \), il existe \( z' \in \mathbb{U} \) tel que \( zz' = 1 \).
\end{exercice}

\begin{exercice}[title=Exercice 9 \quad$\star\star$\quad\ {[}Calculer{]}]{}{}
Résoudre les équations suivantes dans \( \mathbb{C} \) :
\begin{enumerate}
    \item \( z^4 = -1 \)
    \item \( z^3 = i \)
\end{enumerate}
\end{exercice}

\begin{exercice}[title=Exercice 10 \quad$\star\star$\quad\ {[}Calculer{]}]{}{}
Soit \( n \geq 2 \). Résoudre l’équation suivante dans \( \mathbb{C} \) :
\[
z^n = \bar{z}.
\]
\end{exercice}

\begin{exercice}[title=Exercice 11 \quad$\star\star$\quad\ {[}Calculer{]}]{}{}
Soit \( n \geq 2 \). Calculer le produit des racines \( n^{\text{e}} \) de l’unité.
\end{exercice}

\begin{exercice}[title={Exercice 12 \quad$\star\star$\quad\ {[}Calculer, Chercher{]}}]{}{}
Déterminer tous les nombres complexes \( z \) tels que les points d’affixes \( z \), \( z^2 \) et \( z^4 \) soient alignés.
\end{exercice}

\begin{exercice}[title={Exercice 13 \quad$\star\star\star$\quad\ {[}Chercher, Représenter{]}}]{}{}
Soient \( A_1, \dots, A_n \) des points du plan. Peut-on trouver \( n \) points \( B_1, \dots, B_n \) tels que \( A_1, \dots, A_n \) soient les milieux respectifs des segments \([B_1B_2]\), \([B_2B_3]\), \dots, \([B_nB_1]\) ?
\end{exercice}

\begin{exercice}[title={Exercice 14 \quad$\star\star\star$\quad\ {[}Représenter, Calculer, Chercher{]}}]{}{}
Si \( A_1, \dots, A_n \) sont \( n \) points du plan complexe d’affixes respectives \( z_1, \dots, z_n \), on dit qu’un point \( M \) d’affixe \( z \) appartient à l’enveloppe convexe des \((A_k)_{1 \leq k \leq n}\) s’il existe \( \lambda_1, \dots, \lambda_n \in [0;1] \) tels que \( \sum_{k=1}^{n} \lambda_k = 1 \) et tels que :
\[
z = \sum_{k=1}^{n} \lambda_k z_k.
\]
On considère \( P \in \mathbb{C}[X] \). On note \( (a_i)_{1 \leq i \leq s} \) ses racines et, pour tout \( 0 \leq i < s \), \( A_i \) désigne le point d’affixe \( a_i \). Montrer que toutes les racines du polynôme \( P' \) sont contenues dans l’enveloppe convexe des \( (A_i)_{1 \leq i \leq s} \).
\end{exercice}

\begin{exercice}[title={Exercice 15 \quad$\star\star$\quad\ {[}Calculer, Représenter{]}}]{}{}
Soit les nombres complexes :
\( z_1 = \sqrt{2} + i\sqrt{6} \), \( z_2 = 2 + 2i \) et \( Z = \frac{z_1}{z_2} \).
\begin{enumerate}
    \item Écrire \( Z \) sous forme algébrique.
    \item Donner les modules et arguments de \( z_1 \), \( z_2 \) et \( Z \).
    \item En déduire les valeurs de \( \cos \left( \frac{\pi}{12} \right) \) et \( \sin \left( \frac{\pi}{12} \right) \).
    \item Le plan est muni d’un repère orthonormal. On désigne par \( A \), \( B \) et \( C \) les points d’affixes respectives \( z_1 \), \( z_2 \) et \( Z \). Construire à la règle (non graduée) et au compas les points \( A \), \( B \) et \( C \).
    \item Écrire sous forme algébrique le nombre complexe \( Z^{2022} \).
\end{enumerate}
\end{exercice}

\begin{exercice}[title={Exercice 16 \quad$\star\star\star$\quad\ [Représenter]}]{}{}
\begin{enumerate}
    \item Déterminer l’expression de \( \cos(5x) \) en fonction de \( \cos(x) \).
    \item En déduire la valeur de \( \cos \left( \frac{\pi}{5} \right) \), de \( \cos \left( \frac{2\pi}{5} \right) \) et de \( \cos \left( \frac{3\pi}{5} \right) \).
    \item En déduire une construction à la règle (non graduée) et au compas d’un pentagone régulier inscrit dans le cercle trigonométrique.
\end{enumerate}
\end{exercice}

\begin{exercice}[title={Exercice 17 \quad$\star\star\star$\quad\ {[}Représenter, Calculer, Chercher{]}}]
Dans un repère orthonormé, peut-on trouver un triangle équilatéral dont les trois sommets sont à coordonnées entières ?
\end{exercice}

\begin{exercice}[title={Exercice 18 \quad$\star\star\star$\quad\ {[}Représenter, Calculer, Raisonner{]}}]{}{}
Soient \( A \), \( B \) et \( C \) trois points non alignés d’affixe \( a \), \( b \) et \( c \). On note \( j = e^{\frac{2i\pi}{3}} \).
\begin{enumerate}
    \item Montrer que le triangle \( ABC \) est équilatéral direct si, et seulement si,
    \[
    a + bj + cj^2 = 0.
    \]
    \item On ne suppose pas nécessairement que \( ABC \) est équilatéral. On construit à partir de \( ABC \) les trois triangles équilatéraux de base \([AB]\), \([AC]\) et \([BC]\) construits à l’extérieur de \( ABC \). Montrer que les centres de gravité de ces trois triangles forment un triangle équilatéral.
\end{enumerate}
\end{exercice}

\begin{exercice}[title=Exercice 19 \quad$\star\star$\quad\ ]{}{}
Le plan complexe est muni d'un repère orthonormé \( (O ; \vec{u}, \vec{v}) \) et on considère les points \( A \), \( B \) et \( C \) d'affixes respectives :  
\( z_A = \sqrt{2}e^{-\frac{3i\pi}{4}} \), \( z_B = 3 + 3i \) et \( z_C = 7 - i \).

\begin{enumerate}
    \item Déterminer la forme algébrique de \( z_A \) et une forme exponentielle de \( z_B \).
    \item Quelle est la nature du triangle \( ABC \) ?
    \item Démontrer qu'il existe un point \( D \) tel que le quadrilatère \( ABCD \) soit un carré. Déterminer l'affixe du point \( D \).
    \item Déterminer l'affixe du centre du carré \( ABCD \).
    \item On considère l'ensemble des points \( M \) du plan d'affixe \( z \) tels que \( |z - 3 + i| = 4 \). Démontrer que cet ensemble est le cercle circonscrit au triangle \( ABC \).
\end{enumerate}
\end{exercice}


\begin{exercice}[title=Exercice 20 \quad$\star\star$\quad\ ]{}{}
\begin{center}
    \includegraphics[width=0.5\textwidth]{schema_exo68.png}
\end{center}
En électricité, il est souvent nécessaire de réaliser la somme de tensions sinusoïdales de même pulsation \( \omega \).  
On obtient \( v(t) = v_1(t) + v_2(t) \), avec \( v_1(t) = 3\sin\left(\omega t + \frac{\pi}{2}\right) \) et \( v_2(t) = 4\sin(\omega t) \).

\begin{enumerate}
    \item \begin{enumerate}
        \item Déterminer \( |4 + 3i| \).
        \item En utilisant la calculatrice, déterminer un argument de \( 4 + 3i \) à \( 10^{-3} \) près. En déduire une forme exponentielle de \( 4 + 3i \) et de \( 4 - 3i \).
    \end{enumerate}
    \item En utilisant les formules d'Euler, montrer que
       \[
       v(t) = \frac{1}{2i} \left(e^{i\omega t} (4 + 3i) - e^{-i\omega t} (4 - 3i)\right).
       \]
    \item En déduire que \( v(t) = V \sin(\omega t + \phi) \) où \( V \) et \( \phi \) sont deux réels qu'il faudra expliciter.
\end{enumerate}

\end{exercice}

\begin{exercice}[title=Exercice 21 \quad$\star\star$\quad\ ]{}{}
Soit \( z = x + iy \) un nombre complexe, où \( x \) et \( y \) sont des nombres réels.  
À tout nombre complexe \( z \), on associe le nombre noté \( e^z \) défini par \( e^z = e^x (\cos(y) + i\sin(y)) \).

\begin{enumerate}
    \item Vérifier que pour tous nombres complexes \( z_1 \) et \( z_2 \), on a : \( e^{z_1} \times e^{z_2} = e^{z_1 + z_2} \).
    \item Écrire la forme algébrique de \( e^{2 + 3i} \).
    \item Déterminer une solution de l'équation \( e^z = 1 + i\sqrt{3} \). En existe-t-il d'autres ?
\end{enumerate}
\end{exercice}

\begin{exercice}[title=Exercice 22 \quad$\star\star$\quad\ ]{}{}
\begin{center}
    \includegraphics[width=0.5\textwidth]{schema_exo65.png}
\end{center}
En électricité, en régime sinusoïdal utilisant un courant alternatif de pulsation \( \omega \), on utilise les nombres complexes pour représenter les tensions et intensités sinusoïdales.  
On sait, grâce aux lois de l'électricité, que la tension complexe \( U_0 \) aux bornes de la résistance \( R \) est reliée à l'intensité complexe \( I \) par la relation \( U_0 = RI \) et que
\[
U_i = \frac{1}{iC\omega} I + U_0,
\]
où \( C \) est la capacité du condensateur.

On cherche le gain du circuit : \( G = \frac{U_0}{U_i} \).

\begin{enumerate}
    \item Démontrer que
       \[
       G = \frac{iRC\omega}{1 + iRC\omega}.
       \]
    \item Quand \( R = 1000 \) ohms, \( C = 64 \times 10^{-6} \) farads et \( \omega = 60000 \) Hz, démontrer que le module de \( G \) est proche de 1.
    \item Rechercher le sens de « circuit passe-haut ».
\end{enumerate}

\end{exercice}

\begin{exercice}[title=Exercice 23 \quad$\star\star$\quad\ ]{}{}
L'objectif est de déterminer la valeur exacte de
\[
I = \int_0^{\frac{\pi}{2}} \sin^4(x) \, dx.
\]

\begin{enumerate}
    \item Développer \( (a + b)^4 \).
    \item En utilisant une formule d'Euler, prouver que pour tout réel \( x \),
       \[
       \sin^4(x) = \frac{1}{8} ( \cos(4x) - 4\cos(2x) + 3 ).
       \]
    \item Conclure.
\end{enumerate}
\end{exercice}

\begin{exercice}[title=Exercice 24 \quad$\star\star$\quad\ ]{}{}
\begin{enumerate}
    \item En écrivant \( e^{i(\theta + \theta')} \) de deux manières différentes, retrouver les formules d'addition concernant \( \cos(\theta + \theta') \) et \( \sin(\theta + \theta') \).
    \item En déduire que pour tout nombre réel \( \alpha \),
       \[
       \cos^2(\alpha) = \frac{1 + \cos(2\alpha)}{2}
       \]
       et donner une expression de \( \sin(2\alpha) \) en fonction de \( \cos(\alpha) \) et \( \sin(\alpha) \).
    \item Calculer, à l'aide de la question 2, la valeur exacte de \( \cos\left(\frac{\pi}{8}\right) \), puis celle de \( \sin\left(\frac{\pi}{8}\right) \).  
       Vérifier les résultats trouvés à l'aide d'une calculatrice.
    \item Avec une calculatrice, on trouve :
       \[
       \tan\left(\frac{\pi}{12}\right) = 2 - \sqrt{3}.
       \]
       Démontrer ce résultat à l'aide de la question 2.
\end{enumerate}
\end{exercice}

\begin{exercice}[title=Exercice 25 \quad$\star$\quad\ ]{}{}
Le plan complexe est muni d'un repère orthonormé direct \( (O ; \vec{u}, \vec{v}) \). On considère l'application \( f \) qui, à tout point d'affixe \( z \), associe le point \( M' \) d'affixe \( z' \) tel que \( z' = z^2 \).

\begin{enumerate}
    \item Déterminer l'ensemble des points \( M \) qui ont pour image eux-mêmes par \( f \).
    \item En utilisant la forme exponentielle, déterminer la forme exponentielle de :
       \begin{enumerate}
           \item \( f(A) \), où \( A \) est le point d'affixe \( 1 - i \).
           \item \( f(B) \), où \( B \) est le point d'affixe \( 2 + 2\sqrt{3}i \).
           \item \( D \) est le point d'affixe \( \sqrt{2} + i\sqrt{2} \). On cherche les antécédents éventuels de \( D \) par \( f \), c'est-à-dire les points \( M \) tels que \( f(M) = D \).
       \end{enumerate}
    \item En écrivant l'affixe de \( M \) sous la forme \( z = re^{i\theta} \) (\( r > 0 \) et \( \theta \in \mathbb{R} \)), déterminer les valeurs possibles de \( r \) et de \( \theta \).
       \begin{enumerate}
           \item Conclure.
       \end{enumerate}
\end{enumerate}
\end{exercice}

\begin{exercice}[title=Exercice 26 \quad$\star\star\star$\quad\ ]{}{}
On considère un nombre réel \( x \) et un entier naturel non nul \( n \). On cherche à trouver une expression explicite des sommes :  
\( C = 1 + \cos(x) + \cos(2x) + \dots + \cos(nx) \)  
\( S = \sin(x) + \sin(2x) + \dots + \sin(nx) \)  

\begin{enumerate}
    \item Exprimer \( C + iS \) comme une somme d'exponentielles complexes.
    \item En utilisant la somme des termes d'une suite géométrique, en déduire que, pour tout réel \( x \) ne s'écrivant pas sous la forme \( x = 2k\pi \) (k $\in \mathbb{Z}$),
       \[
       C + iS = \frac{1 - e^{i(n+1)x}}{1 - e^{ix}}.
       \]
    \item En factorisant respectivement par \( e^{i\frac{x}{2}} \) et par \( e^{i\frac{(n+1)x}{2}} \), démontrer que :  
       \begin{enumerate}
           \item \( 1 - e^{ix} = -2ie^{i\frac{x}{2}} \sin\left(\frac{x}{2}\right) \),
           \item \( 1 - e^{i(n+1)x} = -2ie^{i\frac{(n+1)x}{2}} \sin\left(\frac{(n+1)x}{2}\right) \).
       \end{enumerate}
    \item Prouver que, pour tout réel \( x \) ne s'écrivant pas sous la forme \( x = 2k\pi \) (k entier relatif), on a
       \[
       C + iS = e^{i\frac{nx}{2}} \times \frac{\sin\left(\frac{(n+1)x}{2}\right)}{\sin\left(\frac{x}{2}\right)}.
       \]
    \item En remarquant que \( C \) est la partie réelle et \( S \) la partie imaginaire de \( C + iS \), en déduire la solution du problème.
    \item Étudier le cas où \( x \) s'écrit sous la forme \( x = 2k\pi \) (k entier relatif).
    \item Applications : \( n \) est un entier naturel supérieur ou égal à 2.
       \begin{enumerate}
           \item Calculer \( \cos\left(\frac{\pi}{n}\right) + \cos\left(\frac{2\pi}{n}\right) + \dots + \cos\left(\frac{(n-1)\pi}{n}\right) \).
           \item Démontrer que
           \(
             \sin\left(\frac{\pi}{n}\right) + \sin\left(\frac{2\pi}{n}\right) + \dots + \sin\left(\frac{(n-1)\pi}{n}\right) = \frac{1}{\tan\left(\frac{\pi}{2n}\right)}.
           \)
       \end{enumerate}
\end{enumerate}
\end{exercice}

\end{multicols}

\end{document}